\documentclass[12pt,a4paper]{article}
\usepackage{task}

\newcommand{\contestName}{MateInfoUB}
\newcommand{\contestDates}{17 Mai 2025}
\newcommand{\contestPlace}{Facultatea de Matematic\v a-Informatic\v a}
\newcommand{\contestRound}{Rund\v a de Test}
%\newcommand{\contestRound}{Rund\v a de Test}
\newcommand{\statementLanguage}{Rom\^ an\v a (Oficial)}
\newcommand{\taskShortName}{A+B}

\setlength{\parskip}{0pt}

\begin{document}
% Tex fragment for task statement headers
% Usage:
%       % Defining parameters required for header
%       \newcommand{\contestName}{Name of the Contest Series}
%       \newcommand{\contestDates}{Month 1 -- 8 20XX}
%       \newcommand{\contestPlace}{City, Country}
%       \newcommand{\contestRound}{Day 1}
%       \newcommand{\statementLanguage}{en (US)}
%       \newcommand{\taskShortName}{task_short_name}
%       % Tex fragment for task statement headers
% Usage:
%       % Defining parameters required for header
%       \newcommand{\contestName}{Name of the Contest Series}
%       \newcommand{\contestDates}{Month 1 -- 8 20XX}
%       \newcommand{\contestPlace}{City, Country}
%       \newcommand{\contestRound}{Day 1}
%       \newcommand{\statementLanguage}{en (US)}
%       \newcommand{\taskShortName}{task_short_name}
%       % Tex fragment for task statement headers
% Usage:
%       % Defining parameters required for header
%       \newcommand{\contestName}{Name of the Contest Series}
%       \newcommand{\contestDates}{Month 1 -- 8 20XX}
%       \newcommand{\contestPlace}{City, Country}
%       \newcommand{\contestRound}{Day 1}
%       \newcommand{\statementLanguage}{en (US)}
%       \newcommand{\taskShortName}{task_short_name}
%       \input{header.tex}
%


%--------------------- tools ----------------------
\makeatletter
% \expandafter for the case that the parameter is given in a command
\newcommand{\escapeUnderscores}[1]{\expandafter\@repl@underscores#1_\relax}
\def\@repl@underscores#1_#2\relax{%
    \ifx \relax #2\relax
        % #2 is empty => finish
        #1%
    \else
        % #2 is not empty => underscore was contained, needs to be replaced
        #1%
        \textunderscore
        % continue replacing
        % #2 ends with an extra underscore so I don't need to add another one
        \@repl@underscores#2\relax
    \fi
}
\makeatother
% -------------------------------------------------

\vspace*{-3em}\hspace*{-.5cm}
\begin{tabular}{ccl}
    \hspace{1mm}\includegraphics[width=2.2cm,valign=b]{logo.jpeg}
    & 
    \begin{minipage}[b]{10cm}
        \setlength{\baselineskip}{1.05\baselineskip}
        \sffamily
        \makebox[0pt][l]{\bfseries \large \contestName}  \\
        \contestDates \\ 
        \contestPlace
    \end{minipage}
    & 
    \begin{minipage}[b]{3.7cm}
        \begin{flushright}
            \makebox[0pt][r]{\ttfamily \bfseries \large \escapeUnderscores{\taskShortName}}  \\[.2em]
            \sffamily
            \contestRound \\
            \statementLanguage
        \end{flushright}
    \end{minipage}
\end{tabular} 
\hrule height .06em
%


%--------------------- tools ----------------------
\makeatletter
% \expandafter for the case that the parameter is given in a command
\newcommand{\escapeUnderscores}[1]{\expandafter\@repl@underscores#1_\relax}
\def\@repl@underscores#1_#2\relax{%
    \ifx \relax #2\relax
        % #2 is empty => finish
        #1%
    \else
        % #2 is not empty => underscore was contained, needs to be replaced
        #1%
        \textunderscore
        % continue replacing
        % #2 ends with an extra underscore so I don't need to add another one
        \@repl@underscores#2\relax
    \fi
}
\makeatother
% -------------------------------------------------

\vspace*{-3em}\hspace*{-.5cm}
\begin{tabular}{ccl}
    \hspace{1mm}\includegraphics[width=2.2cm,valign=b]{logo.jpeg}
    & 
    \begin{minipage}[b]{10cm}
        \setlength{\baselineskip}{1.05\baselineskip}
        \sffamily
        \makebox[0pt][l]{\bfseries \large \contestName}  \\
        \contestDates \\ 
        \contestPlace
    \end{minipage}
    & 
    \begin{minipage}[b]{3.7cm}
        \begin{flushright}
            \makebox[0pt][r]{\ttfamily \bfseries \large \escapeUnderscores{\taskShortName}}  \\[.2em]
            \sffamily
            \contestRound \\
            \statementLanguage
        \end{flushright}
    \end{minipage}
\end{tabular} 
\hrule height .06em
%


%--------------------- tools ----------------------
\makeatletter
% \expandafter for the case that the parameter is given in a command
\newcommand{\escapeUnderscores}[1]{\expandafter\@repl@underscores#1_\relax}
\def\@repl@underscores#1_#2\relax{%
    \ifx \relax #2\relax
        % #2 is empty => finish
        #1%
    \else
        % #2 is not empty => underscore was contained, needs to be replaced
        #1%
        \textunderscore
        % continue replacing
        % #2 ends with an extra underscore so I don't need to add another one
        \@repl@underscores#2\relax
    \fi
}
\makeatother
% -------------------------------------------------

\vspace*{-3em}\hspace*{-.5cm}
\begin{tabular}{ccl}
    \hspace{1mm}\includegraphics[width=2.2cm,valign=b]{logo.jpeg}
    & 
    \begin{minipage}[b]{10cm}
        \setlength{\baselineskip}{1.05\baselineskip}
        \sffamily
        \makebox[0pt][l]{\bfseries \large \contestName}  \\
        \contestDates \\ 
        \contestPlace
    \end{minipage}
    & 
    \begin{minipage}[b]{3.7cm}
        \begin{flushright}
            \makebox[0pt][r]{\ttfamily \bfseries \large \escapeUnderscores{\taskShortName}}  \\[.2em]
            \sffamily
            \contestRound \\
            \statementLanguage
        \end{flushright}
    \end{minipage}
\end{tabular} 
\hrule height .06em

\section*{A plus B}

Se dau două numere $A$ și $B$.

Trebuie să calculați suma lor.

% \subsection*{Detalii de Implementare}

\subsection*{Date de intrare}


Pe prima linie a datelor de intrare se află numerele $A$ și $B$.

\subsection*{Date de ieșire}

Afișați o singură linie care să conțină suma numerelor $A$ și $B$.

\subsection*{Constrângeri}

\begin{itemize}
    \item $1 \leq A \leq 10^{100.000}$
    \item $1 \leq B \leq 10^{100.000}$
\end{itemize}


\subsection*{Subtask-uri}

\begin{enumerate}
    \item ($50$ de puncte) $1 \leq A, B \leq 10^{18}$
    \item ($50$ de puncte) Nicio constrângere suplimentară
\end{enumerate}

\subsection*{Exemplu}

\begin{tabular}{|@{}p{0.5\textwidth}@{}|@{}p{0.5\textwidth}@{}|}
\hline
\multicolumn{1}{|c|}{\bfseries Input Standard (\textit{cin})} &
\multicolumn{1}{c|}{\bfseries Output Standard (\textit{cout})} \\
\hline
\begin{textQuoteCell}
10 5
\end{textQuoteCell} &
\begin{textQuoteCell}
15
\end{textQuoteCell} \\    
\hline
\end{tabular}
\vspace{1em}

% In the first example, $n=5$. The answer is $0$ because...

% In the second example, $n=5$. The answer is $0$ because...

\end{document}